\begin{summarybox}
    \textbf{\textcolor{CSummary}{一句话核心}}

    遇到 $\dfrac{ax+b}{cx+d}$，先把分子凑成“分母的倍数加一个数”。当 $c\ne0$、$ad\ne bc$ 且 $x$ 只受分母不为 $0$ 的限制时，\textbf{函数值恰好只缺 $a/c$ 这一个数}。

    \medskip

    \medskip
    \textbf{\textcolor{CSummary}{先想旧知识，公式就不用死记}}
    \[
      \underbrace{y=\frac3x}_{\text{初中反比例函数}}
      \ \xrightarrow{\ \text{向右移 }1\ }\
      y=\frac3{x-1}
      \ \xrightarrow{\ \text{向上移 }2\ }\
      \boxed{y=2+\frac3{x-1}}.
    \]
    原来碰不到的两条直线 $x=0$、$y=0$，分别变成 $x=1$、$y=2$；
    中心从 $(0,0)$ 变成 $(1,2)$，但\textbf{中心始终不在图像上}。
    更关键的是\textbf{整条 $y=2$ 没有图像点}，所以 $y=2$ 取不到。
    注意：这比仅说“中心那个点取不到”更强，也才真正解释了为什么值域缺一个数。

    \textbf{\textcolor{CSummary}{一般形式只看三件事}}
    \[
       y=\underbrace{\frac ac}_{\text{上下平移}}
       +\frac{\overbrace{(bc-ad)/c^2}^{k\ne0}}{
          x-\underbrace{(-d/c)}_{\text{左右平移}}}.
    \]
    \textbf{先认出 $y=k/x$，再读横向平移、纵向平移，最后写值域。}
    \textbf{$x=-d/c$ 不能代入，$y=a/c$ 不能得到，中心就在 $(-d/c,\,a/c)$。}

    \textbf{\textcolor{CSummary}{30 秒解题步骤}}
    \begin{description}[style=nextline, leftmargin=2.8em, labelsep=0.5em, itemsep=0.35em]
        \item[\textbf{一看条件：}] $c\ne0$、$ad\ne bc$，且没有额外规定 $x$ 的范围。
        \item[\textbf{二拆分式：}] $\displaystyle y=\frac ac+\frac{bc-ad}{c(cx+d)}$。
        \item[\textbf{三写结果：}] 定义域为 $x\ne-d/c$；值域为 $y\ne a/c$；图像中心为 $(-d/c,a/c)$。
    \end{description}

    \medskip
    \textbf{\textcolor{CSummary}{为什么成立（记住最关键的理由）}}

    拆开后剩下的分式不等于 $0$，所以缺少 $a/c$；而对其他每个 $y$ 都能反过来找到合法的 $x$，所以\textbf{恰好只缺这一个值}。

    \medskip
    \textbf{\textcolor{CSummary}{最容易失分的两个地方}}
    \begin{itemize}[leftmargin=*, itemsep=0.5em, topsep=0.4em]
        \item $ad=bc$ 时是常数函数 $y=a/c$（但 $x=-d/c$ 仍不允许），不能再写 $y\ne a/c$。
        \item 若额外要求 $x\ge2$、$x>0$ 等，必须将拆分式和 $x$ 的范围结合起来，重新求值域。
    \end{itemize}

    \textbf{\textcolor{CSummary}{顺手多记一个}}
    \[
        \left(x+\frac dc\right)\left(y-\frac ac\right)
        =\frac{bc-ad}{c^2}.
    \]
    一眼看出它和 $Y=k/X$ 是同一种图像，也就能记住中心 $\left(-d/c,\,a/c\right)$。
\end{summarybox}
